\documentclass[11pt]{article}
\usepackage[margin=1in]{geometry}
\usepackage{amsmath,amssymb,amsthm,mathtools,hyperref}
\newtheorem{theorem}{Theorem}
\newtheorem{lemma}{Lemma}
\newcommand{\tr}{\operatorname{tr}}
\newcommand{\VS}{\operatorname{VS}}
\title{The Vendi--LCR inequality at every order\\MathDB 371345}
\author{Principia Math In Progress}
\date{18 September 2026 (UTC)}
\begin{document}
\maketitle
\noindent\textbf{Scope and prior work.} The current source \cite{N}, version 4,
already proves the comparison at every nonnegative order, and at negative
orders for invertible matrices. The remaining case is a singular matrix at
a negative order. We resolve that case using the orthogonal projection onto
the range, and give a complete proof for all orders. The positive spectrum
convention is the one used in the source's Appendix C.3 code (p.~15).
The audit is by the authoring agent, not an independent peer review.

\section*{Task 1: full comparison, including singular negative orders}
Let $Z=(z_{ij})_{i,j=1}^n$ be real symmetric positive semidefinite, with
$z_{ii}=1$ and $0\leq z_{ij}\leq1$. Let $\lambda_1,\ldots,\lambda_r>0$
be its positive eigenvalues, $r=\operatorname{rank}Z$, and put
\[
 r_i=\sum_j z_{ij},\qquad s_i=\sum_jz_{ij}^2,
 \qquad 1\leq s_i\leq r_i.
\]
For finite $q\ne1$, define
\[
 \VS_q(Z)=\left(n^{-q}\sum_{\ell=1}^r\lambda_\ell^q\right)^{1/(1-q)},
 \qquad
 D_q(Z)=\left(n^{-q}\sum_{i=1}^n r_i^{q-1}\right)^{1/(1-q)}.
\]
Thus zero eigenvalues are omitted, including at nonpositive orders.
At $q=1$ use the continuous limits
\[
 \VS_1=\exp\left(\log n-\frac1n\sum_\ell\lambda_\ell\log\lambda_\ell\right),
 \qquad D_1=\exp\left(\log n-\frac1n\sum_i\log r_i\right).
\]
At the two infinite endpoints use the limits in $q$:
\[
 \VS_{-\infty}=\frac n{\min_\ell\lambda_\ell},\quad
 D_{-\infty}=\frac n{\min_i r_i},\qquad
 \VS_{+\infty}=\frac n{\max_\ell\lambda_\ell},\quad
 D_{+\infty}=\frac n{\max_i r_i}.
\]

\begin{theorem}\label{thm:main}
For every such $Z$ and every $q\in[-\infty,+\infty]$,
$\VS_q(Z)\geq D_q(Z)$.
More precisely, if $P$ is the orthogonal projection onto $\operatorname{ran}Z$
and $h_i=P_{ii}$, then for every finite $q\leq0$,
\begin{equation}\label{eq:chain}
 \sum_{\ell=1}^r\lambda_\ell^q
 \ \geq\ \sum_i h_i^{1-q}
 \ \geq\ \sum_i s_i^{q-1}
 \ \geq\ \sum_i r_i^{q-1}.
\end{equation}
Equality $\VS_q=D_q$ at any finite $q\leq0$ holds exactly when $Z$ is,
up to simultaneous row and column permutation, a direct sum of all-ones
matrices. Such matrices give equality at every order.
\end{theorem}

\begin{proof}[The new case: every nonpositive order and arbitrary rank]
Since $Ze_i$ lies in $\operatorname{ran}Z$,
\[
 1=\langle e_i,Ze_i\rangle=\langle Pe_i,Ze_i\rangle,
 \qquad
 1\leq\|Pe_i\|^2\|Ze_i\|^2=h_i s_i.
\]
In particular $h_i>0$ and $h_i\geq1/s_i\geq1/r_i$.
Choose an orthonormal eigenbasis $v_1,\ldots,v_r$ for the range and write
$w_{i\ell}=|(v_\ell)_i|^2$. Then
\[
 \sum_\ell w_{i\ell}=h_i,\qquad
 \sum_\ell\lambda_\ell w_{i\ell}=z_{ii}=1.
\]
For $q<0$, the function $t\mapsto t^q$ is convex on $(0,\infty)$.
Jensen's inequality with weights $w_{i\ell}/h_i$ gives
\[
 \sum_\ell w_{i\ell}\lambda_\ell^q
 \geq h_i\left(\frac1{h_i}\right)^q=h_i^{1-q}.
\]
Summing over $i$ proves the first inequality in \eqref{eq:chain}, since
$\sum_iw_{i\ell}=1$. For $q=0$ that first inequality is equality:
$r=\sum_i h_i$. The remaining inequalities follow from $h_i\geq1/s_i$
and $s_i\leq r_i$, with $1-q>0$ and $q-1<0$.
Multiplication by $n^{-q}$ and raising to the positive power $1/(1-q)$
proves the score comparison for all finite $q\leq0$.

For equality, the last inequality in \eqref{eq:chain} must be equality
term by term. Hence $s_i=r_i$ for every $i$, and
$\sum_jz_{ij}(1-z_{ij})=0$ forces every entry to be either zero or one.
Represent $Z$ as a Gram matrix of unit vectors. An inner product of one
means two vectors are equal; an inner product of zero means they are
orthogonal. The equal-vector classes therefore give the claimed all-ones
blocks. Conversely, if the block sizes are $b_1,\ldots,b_r$, the positive
eigenvalues are $b_1,\ldots,b_r$ and
\[
 \sum_\ell\lambda_\ell^q=\sum_\ell b_\ell^q
 =\sum_i r_i^{q-1}.
\]
The equality conclusion follows, including the limits at $1$ and infinity.
\end{proof}

\paragraph{Completion at positive orders (already established in \cite{N}).}
For completeness, we give another proof of this known portion.
Set $R=\operatorname{diag}(r_1,\ldots,r_n)$. For every vector $u$,
\[
 u^*(R-Z)u=\frac12\sum_{i,j}z_{ij}|u_i-u_j|^2\geq0,
 \qquad 0\preceq Z\preceq R.
\]
We use the familiar matrix power-order facts
\[
 A\preceq B\Longrightarrow A^\alpha\preceq B^\alpha\quad(0<\alpha\leq1),
 \qquad
 A^{-p}\succeq B^{-p}\quad(0<p\leq1, A\succ0).
\]
They follow from the reversal of order under inversion and the scalar
integrals, applied by spectral calculus,
\[
 x^\alpha=\frac{\sin\pi\alpha}{\pi}
 \int_0^\infty t^{\alpha-1}\frac{x}{x+t}\,dt,
 \qquad
 x^{-p}=\frac{\sin\pi p}{\pi}
 \int_0^\infty\frac{t^{-p}}{x+t}\,dt
\]
for exponents strictly between zero and one; the endpoint cases are direct.
Inversion reverses order because conjugating $A\preceq B$ by $B^{-1/2}$
gives a positive definite matrix bounded above by the identity, whose
inverse is bounded below by the identity.

For $1<q\leq2$, take $\alpha=q-1$. Multiplying the power-order inequality
by $Z$ inside a trace gives
\[
 \tr Z^q\leq\tr(ZR^{q-1})=\sum_i r_i^{q-1}.
\]
For $q\geq2$, use a unit eigenvector $v_\ell$ of $Z$ and scalar Jensen:
\[
 \langle v_\ell,R^{q-1}v_\ell\rangle
 \geq\langle v_\ell,Rv_\ell\rangle^{q-1}
 \geq\lambda_\ell^{q-1}.
\]
Multiplying by $\lambda_\ell$ and summing gives the same trace inequality.
Because $1/(1-q)<0$, it proves $\VS_q\geq D_q$ for $q>1$.

For $0<q<1$, put $p=1-q$. Apply the inverse-power order to
$Z+\varepsilon I\preceq R+\varepsilon I$, multiply by $Z$ inside the
trace, and let $\varepsilon\downarrow0$:
\[
 \sum_{\ell=1}^r\lambda_\ell^q
 =\lim_{\varepsilon\downarrow0}\tr\bigl(Z(Z+\varepsilon I)^{-p}\bigr)
 \geq\tr(ZR^{-p})=\sum_i r_i^{q-1}.
\]
The zero-eigenvalue contributions are identically zero before the limit;
this is a finite-dimensional spectral limit. The positive exponent
$1/(1-q)$ now gives $\VS_q\geq D_q$. Continuity gives $q=1$.
Both spectra and the list of row sums are finite and strictly positive
after removing zero eigenvalues. Taking $q\to-\infty$ and $q\to+\infty$
therefore gives the two endpoint formulas and the corresponding inequalities.
This completes the proof of Theorem \ref{thm:main}.

\paragraph{Why the singular case needs its own argument.}
The full-rank estimate $\VS_q\geq n\geq D_q$ at negative orders does not
apply when $\operatorname{rank}Z<n$. Nor can zero eigenvalues simply be
replaced by small positive numbers: negative powers then diverge, whereas
the intended score sums only over the positive spectrum. Formula
\eqref{eq:chain} works directly on that spectrum, with no such limiting
substitution. If instead one adopts the convention that zero eigenvalues
contribute $+\infty$ at negative order, the comparison is immediate there;
the theorem above proves the stronger, finite, source-intended version.

\section*{Tasks 2--5: proposed follow-up problems}
The following are directions, not results proved in this note.
\begin{description}
\item[Task 2.] Extend the full comparison to a nonuniform probability vector
$p$, comparing the spectrum of $\operatorname{diag}(\sqrt p)Z
\operatorname{diag}(\sqrt p)$ with
$\bigl(\sum_i p_i(Zp)_i^{q-1}\bigr)^{1/(1-q)}$, with appropriate support
conventions.
\item[Task 3.] Classify all equality cases at each finite positive order,
including the Shannon order $q=1$.
\item[Task 4.] At a fixed $q<0$, obtain optimal quantitative stability
bounds relating the trace gap in \eqref{eq:chain} to
$\sum_{i,j}z_{ij}(1-z_{ij})$ and distance from all-ones block matrices.
\item[Task 5.] For prescribed $n$, rank, $q<0$, and a lower bound
$\varepsilon>0$ on the positive eigenvalues, determine the maximum possible
ratio $\VS_q(Z)/D_q(Z)$ and characterize extremizers.
\end{description}

\begin{thebibliography}{1}
\bibitem{N} P.~Nguyen, J.~Couch, R.~Bansal, A.~Morgan, C.~Tam, M.~Li,
Rima Arnaout, and Ramy Arnaout,
\emph{Which Similarity-Sensitive Entropy (Sentropy)?},
arXiv:2511.03849v4, submitted 30 March 2026, PDF dated 1 April 2026.
Theorem 1 and Section III, pp.~3--6; Appendix C.3, p.~15.
\url{https://arxiv.org/abs/2511.03849v4}.
\end{thebibliography}
\end{document}
